\documentclass[12pt,a4paper]{article}

% Codificación y composición del documento.
\usepackage[utf8]{inputenc}
\usepackage[T1]{fontenc}
\usepackage[spanish,es-nodecimaldot]{babel}
\usepackage{lmodern}
\usepackage{geometry}
\geometry{margin=2.5cm}
\usepackage{microtype}
\usepackage{graphicx}
\usepackage{float}
\usepackage{booktabs}
\usepackage{array}
\usepackage{longtable}
\usepackage{enumitem}
\usepackage{xcolor}
\usepackage{listings}
\usepackage{hyperref}
\hypersetup{
    colorlinks=true,
    linkcolor=blue!50!black,
    urlcolor=blue!50!black,
    pdftitle={Memoria tecnica del Administrador de Blog},
    pdfauthor={Bases de Datos Avanzadas}
}

\definecolor{codegray}{RGB}{245,245,245}
\lstset{
    basicstyle=\ttfamily\small,
    backgroundcolor=\color{codegray},
    frame=single,
    rulecolor=\color{black!20},
    breaklines=true,
    columns=fullflexible,
    showstringspaces=false
}

\setlist[itemize]{leftmargin=*}
\setlist[enumerate]{leftmargin=*}

\title{\textbf{Memoria tecnica del Administrador de Blog}\\
\large Proyecto de primera evaluacion\\
\normalsize Bases de Datos Avanzadas}
\author{Aplicacion Python, Oracle Database y PL/SQL}
\date{\today}

\begin{document}

\maketitle
\tableofcontents
\newpage

\section{Resumen ejecutivo}

El proyecto implementa una aplicacion de escritorio para administrar un blog
con una base de datos Oracle. La interfaz esta construida con Python y
CustomTkinter; la persistencia, las restricciones de integridad y las
operaciones principales se resuelven en Oracle mediante PL/SQL.

La solucion se distribuye en tres capas claramente separadas:

\begin{enumerate}
    \item \textbf{Persistencia:} tablas relacionales, claves primarias,
    claves foraneas, relaciones muchos-a-muchos y datos iniciales.
    \item \textbf{Servicios de base de datos:} paquetes PL/SQL que exponen
    operaciones de escritura y consultas con \texttt{SYS\_REFCURSOR}.
    \item \textbf{Aplicacion:} ventana grafica, navegacion, formularios,
    filtros, comentarios y adaptacion de los resultados Oracle para la GUI.
\end{enumerate}

La base de datos se ejecuta dentro de Docker para que el entorno sea
reproducible. El usuario solo necesita iniciar el contenedor, cargar los
scripts SQL y ejecutar la aplicacion Python con las dependencias instaladas.

\section{Relacion con el enunciado}

El enunciado solicita una aplicacion Python que administre un blog en Oracle
mediante procedimientos y funciones almacenadas en PL/SQL. Tambien solicita
un documento PDF con el codigo documentado y evidencias de ejecucion, ademas
del repositorio con el codigo fuente.

La implementacion responde a esas condiciones de la siguiente manera:

\begin{longtable}{p{0.32\textwidth}p{0.58\textwidth}}
\toprule
\textbf{Requisito} & \textbf{Implementacion en el proyecto} \\
\midrule
Aplicacion en Python & \texttt{python\_devops/main\_gui.py} inicia la GUI con CustomTkinter. \\
Base de datos Oracle & Oracle Database Free se ejecuta mediante \texttt{docker-compose.yml}. \\
Persistencia con PL/SQL & \texttt{pl\_sql\_backend/02\_procedures.sql} define los paquetes de usuarios, articulos, comentarios, categorias y etiquetas. \\
Relaciones N:M & \texttt{article\_tags} y \texttt{article\_categories} almacenan las asociaciones. \\
Textos extensos & Articulos y comentarios utilizan columnas \texttt{CLOB}. \\
Codigo documentado & Los modulos Python, SQL y scripts contienen encabezados, docstrings y comentarios tecnicos. \\
Repositorio reproducible & \texttt{docker-compose.yml}, scripts de arranque, dependencias y pruebas permiten reconstruir el entorno. \\
\bottomrule
\end{longtable}

\section{Arquitectura general}

El recorrido completo de una operacion es el siguiente:

\begin{enumerate}
    \item El usuario interactua con una vista de CustomTkinter.
    \item La vista valida los datos basicos del formulario.
    \item La vista llama a una funcion de \texttt{db\_connection.py}.
    \item El adaptador abre una conexion con \texttt{oracledb} y ejecuta un
    procedimiento PL/SQL o una consulta parametrizada.
    \item Oracle aplica restricciones, ejecuta la logica de negocio y confirma
    o revierte la transaccion.
    \item Python transforma el resultado en filas, listas o identificadores.
    \item La vista actualiza los widgets visibles y comunica el resultado al
    usuario.
\end{enumerate}

La conexion no se mantiene abierta entre operaciones. Cada helper utiliza un
gestor de contexto para cerrar cursor y conexion al terminar. Esto evita
conexiones huerfanas y mantiene la responsabilidad transaccional dentro del
procedimiento PL/SQL que realiza la escritura.

\section{Modelo relacional}

\subsection{Entidades principales}

\begin{longtable}{p{0.23\textwidth}p{0.67\textwidth}}
\toprule
\textbf{Tabla} & \textbf{Responsabilidad y columnas relevantes} \\
\midrule
\texttt{users} & Identidad de autores y comentaristas. Incluye \texttt{id}, \texttt{name} y \texttt{email}; el correo es unico. \\
\texttt{articles} & Publicaciones del blog. Incluye \texttt{id}, \texttt{title}, \texttt{pub\_date}, \texttt{text} como CLOB y \texttt{user\_id}. \\
\texttt{comments} & Comentarios asociados a un articulo y a un usuario. Incluye \texttt{content} como CLOB y \texttt{creation\_date}. \\
\texttt{tags} & Etiquetas reutilizables con \texttt{id}, \texttt{name} y \texttt{url}. \\
\texttt{categories} & Categorias reutilizables con \texttt{id}, \texttt{name} y \texttt{url}. \\
\bottomrule
\end{longtable}

Las relaciones de uno a muchos se implementan con claves foraneas:

\begin{itemize}
    \item Un usuario puede publicar muchos articulos.
    \item Un usuario puede escribir muchos comentarios.
    \item Un articulo puede recibir muchos comentarios.
\end{itemize}

Las relaciones muchos-a-muchos se resuelven mediante tablas puente:

\begin{itemize}
    \item \texttt{article\_tags(article\_id, tag\_id)} relaciona articulos y etiquetas.
    \item \texttt{article\_categories(article\_id, category\_id)} relaciona articulos y categorias.
\end{itemize}

Cada tabla puente tiene una clave primaria compuesta que impide registrar dos
veces la misma asociacion. Las claves foraneas utilizan \texttt{ON DELETE
CASCADE}, de modo que una asociacion no queda huerfana al eliminar su articulo
o su taxonomia relacionada.

\subsection{Datos iniciales}

El script \texttt{Schema\_db/01\_schema.sql} carga tres usuarios, cinco
etiquetas, tres categorias y dos articulos. El script primero elimina las
tablas del modelo y posteriormente las crea de nuevo, por lo que su ejecucion
es un reinicio destructivo del volumen de desarrollo.

\section{API PL/SQL}

El archivo \texttt{pl\_sql\_backend/02\_procedures.sql} define cinco paquetes.
Los procedimientos de lectura reciben un parametro de salida
\texttt{SYS\_REFCURSOR}; los adaptadores Python lo convierten en una lista de
filas antes de entregarla a la interfaz.

\begin{longtable}{p{0.25\textwidth}p{0.68\textwidth}}
\toprule
\textbf{Paquete} & \textbf{Procedimientos y funcion} \\
\midrule
\texttt{pkg\_users} & \texttt{insert\_user(p\_name, p\_email)} crea usuarios; \texttt{get\_all\_users(p\_cursor)} devuelve el listado. \\
\texttt{pkg\_articles} & \texttt{create\_article(..., p\_article\_id OUT)} crea el articulo y devuelve su ID; \texttt{assign\_tag} y \texttt{assign\_category} pueblan las tablas puente; \texttt{get\_all\_articles} lista articulos. \\
\texttt{pkg\_comments} & \texttt{add\_comment} crea comentarios; \texttt{get\_by\_article} devuelve el hilo de un articulo. \\
\texttt{pkg\_categories} & \texttt{insert\_category} crea categorias y \texttt{get\_all} las consulta. \\
\texttt{pkg\_tags} & \texttt{insert\_tag} crea etiquetas y \texttt{get\_all} las consulta. \\
\bottomrule
\end{longtable}

Las operaciones de escritura aplican \texttt{COMMIT} cuando terminan
correctamente. Ante un error aplicable, ejecutan \texttt{ROLLBACK} y elevan un
mensaje mediante \texttt{RAISE\_APPLICATION\_ERROR}. Algunos casos definidos
son correo duplicado, usuario inexistente, articulo o taxonomia inexistente y
asociacion repetida.

\subsection{CRUD de taxonomia}

Los paquetes de categorias y etiquetas exponen alta, consulta, edicion y
eliminacion. La interfaz genera el nuevo slug al crear o editar un registro y
utiliza los siguientes procedimientos:

\begin{itemize}
    \item \texttt{pkg\_categories.update\_category} y
    \texttt{pkg\_categories.delete\_category}.
    \item \texttt{pkg\_tags.update\_tag} y
    \texttt{pkg\_tags.delete\_tag}.
\end{itemize}

Cada operacion valida que el identificador exista, confirma la transaccion si
la modifica y revierte los cambios ante un error. Si una taxonomia esta
referenciada por una asociacion, las restricciones de integridad de Oracle
determinan si la eliminacion puede realizarse de acuerdo con el modelo
relacional.

\section{Capa de conexion Python}

\subsection{Configuracion}

Los valores utilizados por Docker y Python son:

\begin{center}
\begin{tabular}{ll}
\toprule
\textbf{Parametro} & \textbf{Valor} \\
\midrule
Host & \texttt{localhost} \\
Puerto & \texttt{1521} \\
Service/PDB & \texttt{FREEPDB2} \\
Usuario & \texttt{blog\_admin} \\
Contrasena de desarrollo & \texttt{admin123} \\
\bottomrule
\end{tabular}
\end{center}

La cadena equivalente es:
\begin{lstlisting}
blog_admin/admin123@localhost:1521/FREEPDB2
\end{lstlisting}

Estas credenciales son exclusivamente locales y no deben utilizarse en
produccion.

\subsection{Conversion de CLOB}

\texttt{db\_connection.py} configura \texttt{outputtypehandler}. Cuando
Oracle devuelve una columna de tipo CLOB, el handler solicita un tipo largo de
texto y entrega un \texttt{str} a Python. Asi, la vista puede mostrar el texto
del articulo o comentario sin invocar manualmente \texttt{.read()}.

\subsection{Adaptadores disponibles}

\begin{itemize}
    \item \texttt{get\_connection}: abre una conexion configurada.
    \item \texttt{call\_procedure}: ejecuta procedimientos sin salida.
    \item \texttt{call\_procedure\_returning\_id}: registra el parametro OUT
    y devuelve el ID generado.
    \item \texttt{call\_procedure\_with\_cursor}: consume un
    \texttt{SYS\_REFCURSOR} y devuelve filas Python.
    \item \texttt{fetch\_article\_text}: obtiene el cuerpo completo de un articulo.
    \item \texttt{fetch\_article\_taxonomy}: devuelve nombres de etiquetas y
    categorias asociadas a un articulo.
\end{itemize}

\section{Interfaz grafica y flujos de uso}

\subsection{Ventana principal}

\texttt{main\_gui.py} crea \texttt{BlogAdminApp}, configura el modo oscuro y
monta una barra lateral con las vistas Inicio, Usuarios, Categorias y
Etiquetas. La zona central se limpia y reconstruye al cambiar de seccion.
El boton de publicacion abre una ventana independiente para no perder el
contexto de la ventana principal.

\subsection{Feed de articulos}

\texttt{FeedView} obtiene los articulos, busca el nombre del autor y consulta
la taxonomia de cada articulo. Cada resultado se representa con
\texttt{ArticleCard}, que muestra titulo, autor, fecha y chips de categorias y
etiquetas. Los controles permiten ordenar por fecha y filtrar por una
categoria o etiqueta.

\texttt{utils.safe\_get\_all} intenta primero el procedimiento PL/SQL. Si la
lectura no devuelve filas o falla, utiliza una consulta SQL de respaldo para
los listados conocidos. Esta decision mantiene operativa la interfaz sin
confundir un error de lectura con una escritura exitosa.

\subsection{Detalle y comentarios}

Al seleccionar una tarjeta, \texttt{ArticleDetailView} recibe el diccionario
del articulo. Si el texto no esta cargado, llama a \texttt{fetch\_article\_text}.
Despues consulta \texttt{pkg\_comments.get\_by\_article}, muestra cada
comentario con autor y fecha, y ofrece un formulario con usuario y contenido.

Al pulsar \textit{Comentar}, la interfaz valida los campos y ejecuta solamente
\texttt{pkg\_comments.add\_comment}. La confirmacion de la transaccion es
responsabilidad del procedimiento PL/SQL; no se abre una segunda conexion para
realizar un \texttt{commit} independiente.

\subsection{Publicacion de articulos}

\texttt{NewArticleView} solicita titulo, texto y autor. Las etiquetas y
categorias se cargan como listas de casillas de seleccion multiple. El flujo
es:

\begin{enumerate}
    \item Validar titulo, texto y usuario.
    \item Ejecutar \texttt{pkg\_articles.create\_article}.
    \item Recibir el ID mediante el parametro OUT.
    \item Ejecutar \texttt{pkg\_articles.assign\_tag} por cada etiqueta marcada.
    \item Ejecutar \texttt{pkg\_articles.assign\_category} por cada categoria marcada.
    \item Cerrar la ventana y recargar el feed.
\end{enumerate}

\subsection{Usuarios, categorias y etiquetas}

\texttt{UsersView} consulta usuarios y registra nuevos mediante
\texttt{pkg\_users.insert\_user}. Las vistas de taxonomia comparten
\texttt{\_TaxonomyPanel}: generan automaticamente un slug URL, llaman al
procedimiento de alta y recargan la lista. Los listados se presentan dentro de
contenedores desplazables para soportar una cantidad variable de registros.

\section{Inicializacion del entorno}

\subsection{Arranque automatizado}

Desde la raiz del repositorio se ejecuta:

\begin{lstlisting}[language=bash]
chmod +x start.sh install_python.sh
./start.sh
\end{lstlisting}

El script comprueba Docker, levanta el servicio \texttt{blog\_oracle\_db},
espera su disponibilidad, verifica si existen tablas, carga el esquema y los
paquetes cuando es necesario, crea el entorno virtual y ejecuta la prueba de
conexion.

\subsection{Arranque manual}

\begin{enumerate}
    \item Iniciar Oracle: \texttt{docker compose up -d}.
    \item Esperar el estado \texttt{healthy} con \texttt{docker compose ps}.
    \item Cargar \texttt{Schema\_db/01\_schema.sql} y
    \texttt{pl\_sql\_backend/02\_procedures.sql} si el volumen es nuevo.
    \item Crear \texttt{python\_devops/.venv} e instalar
    \texttt{requirements.txt}.
    \item Ejecutar \texttt{test\_connection.py}.
    \item Abrir la GUI con \texttt{.venv/bin/python main\_gui.py}.
\end{enumerate}

En macOS se necesita Tcl/Tk 8.6 o superior. Si el Python del sistema utiliza
Tk 8.5, \texttt{install\_python.sh} descarga un Python 3.12 compatible y
recrea el entorno virtual.

\section{Pruebas y evidencias}

\subsection{Pruebas automatizadas disponibles}

\begin{itemize}
    \item \texttt{test\_connection.py}: verifica la conexion, una consulta
    basica y el acceso a Oracle.
    \item \texttt{test\_integration.py}: comprueba conexion, CLOB, datos
    semilla e invocacion de los cinco paquetes PL/SQL.
    \item \texttt{py\_compile}: valida la sintaxis de todos los modulos
    Python.
    \item \texttt{docker compose config}: valida la configuracion del servicio
    Oracle.
\end{itemize}

\texttt{test\_integration.py} escribe datos de prueba y espera las cantidades
del seed inicial. Por ese motivo no es idempotente: al ejecutarlo varias veces
sobre el mismo volumen pueden aparecer conteos superiores o errores por
duplicados. Para una ejecucion limpia se utiliza
\texttt{docker compose down -v}, seguido de una inicializacion completa.

\subsection{Evidencias visuales para el PDF}

El producto solicitado incluye capturas de la ejecucion. Para completar la
entrega visual se recomienda incorporar, despues de ejecutar la aplicacion,
las siguientes evidencias:

\begin{enumerate}
    \item Ventana principal con el feed y los filtros.
    \item Detalle de un articulo con su texto y comentarios.
    \item Formulario de publicacion con varias etiquetas y categorias marcadas.
    \item Vista de usuarios con el formulario de alta.
    \item Vista de categorias y etiquetas con sus slugs.
    \item Terminal con \texttt{test\_connection.py} y la prueba de integracion.
    \item Terminal con \texttt{docker compose ps} mostrando Oracle saludable.
\end{enumerate}

Las capturas deben mostrar la aplicacion funcionando, no unicamente el codigo.
Se pueden insertar en esta seccion con \texttt{graphicx} usando, por ejemplo:

\begin{lstlisting}[language={[LaTeX]TeX}]
\begin{figure}[H]
    \centering
    \includegraphics[width=0.9\textwidth]{evidencias/feed.png}
    \caption{Feed de articulos con filtros y taxonomia.}
\end{figure}
\end{lstlisting}

\section{Estructura del repositorio}

\begin{longtable}{p{0.35\textwidth}p{0.53\textwidth}}
\toprule
\textbf{Ruta} & \textbf{Contenido} \\
\midrule
\texttt{docker-compose.yml} & Servicio Oracle Database Free, volumen persistente y healthcheck. \\
\texttt{Schema\_db/01\_schema.sql} & Tablas, restricciones, relaciones y datos iniciales. \\
\texttt{pl\_sql\_backend/02\_procedures.sql} & Paquetes y procedimientos PL/SQL. \\
\texttt{python\_devops/db\_connection.py} & Conexiones, cursores, CLOB y adaptadores PL/SQL. \\
\texttt{python\_devops/main\_gui.py} & Ventana raiz y navegacion. \\
\texttt{python\_devops/views/} & Feed, detalle, publicacion, usuarios y taxonomia. \\
\texttt{python\_devops/widgets/} & Tarjetas, chips y componentes reutilizables. \\
\texttt{python\_devops/test\_*.py} & Pruebas de conexion e integracion. \\
\texttt{start.sh} & Preparacion automatizada del entorno. \\
\texttt{install\_python.sh} & Instalacion de Python con Tcl/Tk compatible en macOS. \\
\bottomrule
\end{longtable}

\section{Conclusiones}

La solucion separa de forma explicita presentacion, acceso a datos y logica
de base de datos. Oracle conserva la integridad referencial y la logica de
escritura; Python administra la interaccion, la validacion de entrada y la
representacion de resultados. Docker permite reproducir la base de datos y los
scripts de prueba verifican que las piezas puedan comunicarse.

El flujo de publicacion demuestra el requisito N:M al crear un articulo,
recibir su ID y asignarle multiples etiquetas y categorias. El flujo de
comentarios demuestra la relacion entre usuarios, articulos y comentarios. La
documentacion del codigo y las evidencias visuales completan la entrega
solicitada por el proyecto.

\end{document}
